\documentclass{amsart}
% For dealing with references we use the comment environment
\usepackage{verbatim}
\newenvironment{reference}{\comment}{\endcomment}
%\newenvironment{reference}{}{}
\newenvironment{slogan}{\comment}{\endcomment}
\newenvironment{history}{\comment}{\endcomment}

% For commutative diagrams we use Xy-pic
\usepackage[all]{xy}

% We use 2cell for 2-commutative diagrams.
\xyoption{2cell}
\UseAllTwocells

% We use multicol for the list of chapters between chapters
\usepackage{multicol}

% This is generally recommended for better output
\usepackage{lmodern}
\usepackage[T1]{fontenc}

% For cross-file-references

% Package for hypertext links:
\usepackage{hyperref}

% For any local file, say "hello.tex" you want to link to please

% Theorem environments.
%
\theoremstyle{plain}
\newtheorem{theorem}[subsection]{Theorem}
\newtheorem{proposition}[subsection]{Proposition}
\newtheorem{lemma}[subsection]{Lemma}

\theoremstyle{definition}
\newtheorem{definition}[subsection]{Definition}
\newtheorem{example}[subsection]{Example}
\newtheorem{exercise}[subsection]{Exercise}
\newtheorem{situation}[subsection]{Situation}

\theoremstyle{remark}
\newtheorem{remark}[subsection]{Remark}
\newtheorem{remarks}[subsection]{Remarks}

\numberwithin{equation}{subsection}

% Macros
%
\def\lim{\mathop{\mathrm{lim}}\nolimits}
\def\colim{\mathop{\mathrm{colim}}\nolimits}
\def\Spec{\mathop{\mathrm{Spec}}}
\def\Hom{\mathop{\mathrm{Hom}}\nolimits}
\def\Ext{\mathop{\mathrm{Ext}}\nolimits}
\def\SheafHom{\mathop{\mathcal{H}\!\mathit{om}}\nolimits}
\def\SheafExt{\mathop{\mathcal{E}\!\mathit{xt}}\nolimits}
\def\Sch{\mathit{Sch}}
\def\Mor{\mathop{\mathrm{Mor}}\nolimits}
\def\Ob{\mathop{\mathrm{Ob}}\nolimits}
\def\Sh{\mathop{\mathit{Sh}}\nolimits}
\def\NL{\mathop{N\!L}\nolimits}
\def\CH{\mathop{\mathrm{CH}}\nolimits}
\def\proetale{{pro\text{-}\acute{e}tale}}
\def\etale{{\acute{e}tale}}
\def\QCoh{\mathit{QCoh}}
\def\Ker{\mathop{\mathrm{Ker}}}
\def\Im{\mathop{\mathrm{Im}}}
\def\Coker{\mathop{\mathrm{Coker}}}
\def\Coim{\mathop{\mathrm{Coim}}}

% Boxtimes
%
\DeclareMathSymbol{\boxtimes}{\mathbin}{AMSa}{"02}

%
% Macros for moduli stacks/spaces
%
\def\QCohstack{\mathcal{QC}\!\mathit{oh}}
\def\Cohstack{\mathcal{C}\!\mathit{oh}}
\def\Spacesstack{\mathcal{S}\!\mathit{paces}}
\def\Quotfunctor{\mathrm{Quot}}
\def\Hilbfunctor{\mathrm{Hilb}}
\def\Curvesstack{\mathcal{C}\!\mathit{urves}}
\def\Polarizedstack{\mathcal{P}\!\mathit{olarized}}
\def\Complexesstack{\mathcal{C}\!\mathit{omplexes}}
% \Pic is the operator that assigns to X its picard group, usage \Pic(X)
% \Picardstack_{X/B} denotes the Picard stack of X over B
% \Picardfunctor_{X/B} denotes the Picard functor of X over B
\def\Pic{\mathop{\mathrm{Pic}}\nolimits}
\def\Picardstack{\mathcal{P}\!\mathit{ic}}
\def\Picardfunctor{\mathrm{Pic}}
\def\Deformationcategory{\mathcal{D}\!\mathit{ef}}

\setlength{\emergencystretch}{3em}
\begin{document}
\title[Stacks Project, Tag 07VE]{304 --- Stability of exact complexes under Artin-Rees controlled perturbation}
\maketitle
\noindent Copyright (C) 2005--2025 Johan de Jong. Stacks Project authors. Source: \href{https://stacks.math.columbia.edu/tag/07VE}{Tag 07VE}.

\noindent Editorial difficulty note (not author text). Difficulty H1: The proof iteratively corrects a vector by the perturbed image while raising its I-adic order and preserving its image. It then proves a uniform Artin-Rees bound for the perturbed map and closes exactness using Krull intersection; these quantitative coupled corrections exceed a routine application of Artin-Rees..

\noindent Editorial provenance note (not author text). History: excerpt from revision \texttt{a0444\allowbreak 6e57e\allowbreak c1fbc\allowbreak 252a8\allowbreak 71afc\allowbreak ec775\allowbreak 2fb28\allowbreak 07b14}, more-algebra.tex, lines 194--245. Reference presentation was adapted; original-author context and core support blocks were retained or added. The mathematical wording is unchanged. Licensed under GNU FDL 1.2 or later; no invariant sections or cover texts. See ../licenses/COPYING. Context blocks reproduce author definitions and supporting results; they are not additional counted problems.

Let $A$ be a Noetherian ring and let $I \subset A$ be an ideal. Given a
homomorphism $f : M \to N$ of finite $A$-modules there exists a $c \geq 0$
such that
$$
f(M) \cap I^nN \subset f(I^{n - c}M)
$$
for all $n \geq c$, see Algebra, Lemma \href{https://stacks.math.columbia.edu/tag/00IO}{lemma map AR (Tag 00IO)}. In this
situation we will say {\it $c$ works for $f$ in the Artin-Rees lemma}.



\begin{lemma}
\label{lemma-approximate-complex}
Let $A$ be a Noetherian ring. Let $I \subset A$ be an ideal contained in
the Jacobson radical of $A$. Let
$$
S : L \xrightarrow{f} M \xrightarrow{g} N
\quad\text{and}\quad
S' : L \xrightarrow{f'} M \xrightarrow{g'} N
$$
be two complexes of finite $A$-modules as shown. Assume that
\begin{enumerate}
\item $c$ works in the Artin-Rees lemma for $f$ and $g$,
\item the complex $S$ is exact, and
\item $f' = f \bmod I^{c + 1}M$ and $g' = g \bmod I^{c + 1}N$.
\end{enumerate}
Then $c$ works in the Artin-Rees lemma for $g'$ and the
complex $S'$ is exact.
\end{lemma}

\begin{proof}
We first show that $g'(M) \cap I^nN \subset g'(I^{n - c}M)$ for $n \geq c$.
Let $a$ be an element of $M$ such that $g'(a) \in I^nN$. We want to
adjust $a$ by an element of $f'(L)$, i.e, without changing $g'(a)$, so
that $a \in I^{n-c}M$. Assume that $a \in I^rM$, where $r < n - c$.
Then
$$
g(a) = g'(a) + (g - g')(a) \in
I^n N + I^{r + c + 1}N = I^{r + c + 1}N.
$$
By Artin-Rees for $g$ we have $g(a) \in g(I^{r + 1}M)$. Say $g(a) = g(a_1)$
with $a_1 \in I^{r + 1}M$. Since the sequence $S$ is exact, $a - a_1 \in f(L)$.
Accordingly, we write $a = f(b) + a_1$ for some $b \in L$.
Then $f(b) = a - a_1 \in I^rM$. Artin-Rees for $f$ shows that
if $r \geq c$, we may replace $b$ by an element of $I^{r - c}L$.
Then in all cases, $a = f'(b) + a_2$, where
$a_2 = (f - f')(b) + a_1 \in I^{r + 1}M$. (Namely, either $c \geq r$
and $(f - f')(b) \in I^{r + 1}M$ by assumption, or $c < r$ and
$b \in I^{r - c}$, whence again $(f - f')(b) \in I^{c + 1} I^{r - c} M =
I^{r + 1}M$.) So we can adjust $a$ by the element $f'(b) \in f'(L)$ to
increase $r$ by $1$.

\medskip\noindent
In fact, the argument above shows that
$(g')^{-1}(I^nN) \subset f'(L) + I^{n - c}M$ for all $n \geq c$.
Hence $S'$ is exact because
$$
(g')^{-1}(0) = (g')^{-1}(\bigcap I^nN) \subset
\bigcap f'(L) + I^{n - c}M = f'(L)
$$
as $I$ is contained in the Jacobson radical of $A$, see Algebra, Lemma
\href{https://stacks.math.columbia.edu/tag/00IQ}{lemma intersection powers ideal module (Tag 00IQ)}.
\end{proof}
\end{document}
